\section{The texts in their biblical and liturgical contexts}

\subsection{The Introit: peace promised to a people who must wait}

The Introit combines portions of Jeremiah 29:11,12,14. Its speaker is the Lord, and its subject is his intention toward the displaced people: peace, a hearing for their prayer, and gathering from captivity. Jeremiah's full letter makes that promise demanding. It instructs the exiles to build houses, plant gardens, form households and seek the peace of the city in which they live. They must conduct a life there while awaiting the restoration God promises. False prophets offer an earlier escape. The letter requires patient fidelity in circumstances the recipients would prefer to leave.\footnote{Jeremiah 29:1--14, especially 5--14.}

Jerome reads the promise against that false expectation of an immediate return. He also notices the relation between God's action and human prayer. God has power to fulfill his promise without their assistance, yet he summons them to call upon him and seek him with the whole heart. Prayer is therefore part of the promised restoration's human response. A divine promise does not make its recipients passive.\footnote{Jerome, \textit{Commentary on Jeremiah}, book V, on Jeremiah 29:10--14, PL 24, cols.~893--894.}

The accompanying verse, Psalm 84(85):2, remembers blessing and the reversal of Jacob's captivity. The complete psalm proceeds from remembered favor to fresh petitions for mercy. It allows gratitude for deliverance already received and need for further help to be spoken together. Augustine reads the release from captivity through forgiveness: the next verse's remission of sins explains the spiritual bondage from which Christ frees his people. He also treats the prayer for conversion later in the psalm as an acknowledgment that conversion itself comes from God.\footnote{Psalm 84(85):2--8; Augustine, \textit{Expositions of the Psalms}, 84.2--6, NPNF heading Psalm 85.} That interpretation gives the entering assembly a way to confess its own need without erasing the letter's concrete concern for the exiles.

\subsection{The Collect: thought, speech and deed}

The Collect contains a close connection between meditation and action. Its object, \textit{rationabilia}, is rendered in the historical English as ``what is reasonable and just.'' The prayer asks that sustained attention to such things become conduct acceptable to God. Speech is included alongside deeds. A person's understanding may become visible in an honest answer, a restrained judgment or the fulfillment of a responsibility as readily as in an explicitly religious explanation.

Schuster's commentary on this Collect keeps that connection intact: attention is directed toward reasonable things so that words and works please God.\footnote{Ildefonso Schuster, \textit{The Sacramentary}, trans. Arthur Levelis-Marke, vol.~I (1924), ``Sixth Sunday after the Epiphany,'' p.~416.} The prayer's form matters as well as its object. Worshippers ask God to grant this integrity. The ability to think well and act accordingly is sought as a gift. Later, Jerome's account of reason and desire will provide a more detailed explanation of the human capacities involved; he is expounding the Gospel, while Schuster comments directly on this prayer.

\subsection{The Epistle: thanksgiving for a community's conversion}

The appointed lesson begins with thanksgiving and ends with expectation of Jesus from heaven. Between them Paul recalls faith's work, charity and endurance, the manner of his preaching, the believers' reception of it, and the report of their conversion. He thanks God for people whose acts he can actually name. The thanksgiving neither removes their agency nor leaves them to claim that they produced their conversion unaided.

Two transitions determine the passage's argument. First, the recipients become imitators of Paul and of the Lord through their reception of the word in affliction. Second, those imitators themselves become an example to believers in Macedonia and Achaia. The word has reached them and its effects have become known beyond them. The report concerns a changed allegiance: they have turned from idols to serve the living and true God. Its final object is the Son whom God raised from the dead, who delivers from the coming wrath.\footnote{1 Thessalonians 1:2--10.}

The rest of the letter gives this conduct further substance. Paul recalls his care for them and his concern during separation; he urges holiness, brotherly charity and responsible work; he teaches concerning resurrection and vigilance. The opening's hope is therefore already connected with the life to which the whole letter exhorts its recipients.\footnote{1 Thessalonians 2--5.} Chrysostom and Aquinas develop the opening in different grammatical detail. Chrysostom explains the labor involved in charity, especially when attachment to fellow believers becomes costly. Aquinas joins both work and labor to faith, then treats the abundance of charitable works and hope's power to sustain endurance. Their shared conclusion is that the virtues have effects in action.\footnote{John Chrysostom, \textit{Homilies on First Thessalonians}, homily 1, on 1:2--7; Thomas Aquinas, \textit{Commentary on First Thessalonians}, chapter 1, lecture 1 [87862].}

\subsection{The Gradual: praise within a psalm of distress}

The Gradual selects the declaration of rescue and the promise of praise from Psalm 43(44):8--9. In the full psalm, the people first recall God's help to their fathers and confess that possession of the land was his gift. They refuse reliance on their own weapons. Immediately after the selected verses, however, their language changes to complaint: they now suffer defeat and humiliation. They maintain that they have not abandoned God's covenant.\footnote{Psalm 43(44):2--10,18--23.}

That context matters when the Gradual follows an Epistle which expressly mentions tribulation. Confidence and suffering are present within each biblical context. Bellarmine explains the rescue as the memory of actual benefits: the speaker acknowledges that past escapes were due to God. He then follows the psalm into its present lament. Augustine gives the past tense a prophetic interpretation as well. Future deliverance can be spoken of as accomplished because it is certain with God; everlasting praise looks toward the promised kingdom.\footnote{Robert Bellarmine, \textit{Commentary on the Book of Psalms}, Psalm 43, explanations 7--9; Augustine, \textit{Expositions of the Psalms}, 43.6--8, NPNF heading Psalm 44.} These accounts differ about the immediate temporal force of the rescue. Both leave room for present affliction in a person who trusts God.

\subsection{The Alleluia: calling from the depths}

The Alleluia uses the beginning of Psalm 129(130), the \textit{De profundis}. The Latin title means ``Out of the depths.'' The selected words ask to be heard. The full psalm goes on to name iniquity, forgiveness, waiting and redemption. Its opening can therefore be heard as a cry from a person who knows both misery and guilt, without requiring every suffering person to explain a particular calamity as punishment for a particular sin.\footnote{Psalm 129(130):1--8.}

Augustine begins with Jonah praying in the sea creature's belly. Physical confinement cannot prevent a prayer from reaching God. He then asks what depth the congregation inhabits and describes mortal life, burdened by sin. The most desperate condition is the refusal to cry at all: prosperity can conceal a person's need so thoroughly that he laughs at the suggestion of prayer. The cry already acknowledges a possible deliverer.\footnote{Augustine, \textit{Expositions of the Psalms}, 129.1--2, NPNF heading Psalm 130.} Placed before the Gospel, this petition lets a needy assembly ask for the hearing that God promises in the Introit.

\subsection{The Gospel: two images and a prophetic citation}

Matthew 13:31--35 belongs to the discourse of parables. Jesus has already told of seed failing in several kinds of ground and of weeds growing among wheat. The mustard seed and leaven describe growth despite a beginning whose size gives little indication of the result. The seed grows into a plant with branches where birds lodge. The leaven is concealed within a large quantity of meal until it has affected the whole. Each comparison asks the hearer to follow a change through time.\footnote{Matthew 13:1--35, especially 18--33.}

The ending turns from the images to the manner of teaching. Matthew cites Psalm 77(78):2, whose speaker promises parables and then recounts Israel's history. Jerome identifies Asaph and explains that this history contains mysteries to be understood. Chrysostom says Jesus' parables invite inquiry; the statement that he spoke without a parable does not describe every utterance of his ministry. It concerns this occasion. The disciples' subsequent request for an explanation shows the kind of attention such teaching requires.\footnote{Psalm 77(78):1--8; Matthew 13:34--36; Jerome, \textit{Commentary on Matthew}, II, on 13:34--35, PL 26, cols.~92--93; Chrysostom, \textit{Homilies on Matthew}, 47.1.} The interpretations below develop the two Fathers' different explanations of the seed and the three measures rather than assigning a single meaning to every feature.

\subsection{The Offertory: the same petition while gifts are brought}

The Offertory returns to the beginning of Psalm 129(130), but its form differs from the Alleluia: the opening cry is repeated after the request for a hearing. That repetition belongs to this chant; it is not an additional verse in the psalm. The assembly has heard Christ's teaching and now again speaks as a people needing mercy.

Schuster relates this Offertory to the humility of the gifts brought to God. They have already been received from his goodness, and their presentation is accompanied by an appeal to his mercy.\footnote{Schuster, \textit{The Sacramentary}, vol.~III (1927), p.~188, commentary on the Offertory \textit{De profundis}.} As the worshipper presents gifts received from God, he asks for mercy.

\subsection{The Secret: four effects sought for the people}

The Secret is exceptionally brief. Its four verbs ask that the oblation cleanse, renew, govern and protect ``us.'' Cleansing concerns what needs to be removed; renewal concerns the life to be restored; governance concerns its direction; protection concerns its preservation. Their sequence permits a fuller account of change than the single idea of relief from trouble. Schuster expressly draws attention to these four effects of the sacrifice.\footnote{Schuster, \textit{The Sacramentary}, I, p.~417.}

The prayer names the offering and asks God to bring about effects in the people. It therefore relates the sacrifice to their continued life. A community whose faith becomes known through conduct needs precisely such cleansing and direction. A person whose understanding has been awakened still needs protection against a return to former habits. Those connections will be developed separately in the two interpretations.

\subsection{The Preface: the God who receives thanksgiving}

The appointed Preface addresses the Father and confesses him with his only-begotten Son and the Holy Spirit as one God and Lord. It distinguishes the persons while confessing unity of essence and equality of majesty. The phrase concerning revelation explains the source of this confession: what is believed about the Father's glory is confessed of the Son and the Spirit without inequality of Godhead. The thanksgiving then names the angels' unceasing praise.\footnote{\textit{Missale Romanum} (1962), p.~293, no.~1082.}

This confession gives definite content to the Epistle's service of the living and true God. Paul names God the Father, the Lord Jesus Christ and the Holy Spirit within the appointed lesson. The Preface addresses that divine life in worship. Its distinction of persons is also relevant when Jerome's interpretation of the Gospel reaches the number three. A likeness involving three quantities cannot by itself establish what the Church confesses about God; Jerome himself warns against that use of the parable.

\subsection{The Communion: confidence and forgiveness}

The Communion takes its assurance concerning prayer from Mark 11:24. In Mark's narrative, Jesus has acted in the Temple and the disciples have seen the fig tree withered. He teaches them faith in God and then directs them to forgive when they stand praying. The short chant selects the assurance; the surrounding instruction describes the praying person.\footnote{Mark 11:12--26.}

Bede comments on the sequence of faith, the mountain saying and forgiveness. He speaks of faith working through charity. Those who have not attained the perfection associated with extraordinary works are to ask for pardon and eternal life, while forgiving others who have sinned against them.\footnote{Bede, \textit{Commentary on Mark}, on Mark 11:22--25, paragraphs beginning \textit{Et respondens Jesus} and \textit{Et cum stabitis ad orandum}.} Thus the assurance invites real trust, and its application includes an obligation toward the neighbor. The desire to receive from God is examined by the willingness to forgive.

\subsection{The Postcommunion: an appetite that remains after nourishment}

The final proper prayer speaks as one already fed and asks for continuing desire for the same heavenly gifts. Ordinary eating usually reduces hunger. Here reception is followed by a request that desire endure, because the gifts sustain true life. Schuster singles out this feature when he comments on the prayer.\footnote{Schuster, \textit{The Sacramentary}, I, p.~417.} The historical English's ``as they preserve our life'' is its own wording; the Latin's concern with living truly remains available for explanation beside it.

The repeated desire is for the gifts received, so it has a definite object. It concerns continued dependence on what nourishes communion with God. For the people described in the Epistle, that dependence extends into service, endurance and expectation of the Son. For an individual being formed by the Gospel, it involves learning to want what enables that life. The Mass ends its proper prayers by asking that the reception just accomplished bear fruit in a desire that continues.
